\section{Noyaux généalogiques de chemins : composition finie, réalisation unitaire et limite cylindrique}
\label{sec:integral-genealogica-caminos-rev2}
\index{Feynman!intégrale généalogique de chemins}
\index{chemin!catégorie enrichie}
\index{intégrale!cylindrique}

La composition matricielle démontrée dans cette section établit exactement que chaque entrée de
\(U^R\) est la somme des poids ordonnés des routes de longueur \(R\).
L’étape suivante conserve l’information qu’une route visible n’exprime pas et
organise les raffinements finis dans un espace compact de généalogies. De
cette manière, on distingue trois résultats : l’identité combinatoire finie,
l’intégrale cylindrique sur la limite des chemins et la réalisation de cette
intégrale comme propagateur d’un système physique concret.

\subsection{Catégorie de routes et poids fonctoriel
}

Soit \(\mathsf P_{\rm TPK}\) la catégorie libre engendrée par les transitions
admissibles du TPK. Ses objets sont des états enrichis de frontière et ses
morphismes sont des suites composables de transitions. La composition est la
concaténation et le chemin vide est l’identité. Soit \(\mathsf P_X\)
la catégorie des routes visibles obtenue en appliquant le lecteur aux
objets et aux transitions ; ses compositions sont les concaténations visibles.
L’application
\[
 q_{\rm vis}:\mathsf P_{\rm TPK}\longrightarrow\mathsf P_X
\]
exprime la route observable et conserve dans sa fibre les différences de feuillet,
d’orientation, de retenue, de bloc, de frontière et de mémoire.

\begin{definition}[Poids généalogique
]
\label{def:peso-genealogico-caminos-rev2}
Soit \(\mathcal B\) une algèbre réelle normée, éventuellement non commutative, et
soit \(\mathbf B_{\mathcal B}\) la catégorie à un seul objet \(\ast\) dont le
monoïde des endomorphismes est
\(\operatorname{End}(\ast)=(\mathcal B,\cdot,1)\). Un poids généalogique est un foncteur

\[
 W:\mathsf P_{\rm TPK}\longrightarrow\mathbf B_{\mathcal B}
\]
tel que
\[
 W(\operatorname{id}_x)=1,
 \qquad
 W(\gamma_2\circ\gamma_1)=W(\gamma_2)W(\gamma_1).
\]
L’ordre du produit coïncide avec l’ordre causal de la route.

\end{definition}

Pour des états enrichis \(\widetilde i,\widetilde f\), le noyau à la
profondeur \(R\) est
\begin{equation}
 \mathcal K_R(\widetilde f,\widetilde i)
 =\sum_{\substack{
   \gamma\in\mathsf P_{\rm TPK}(\widetilde i,\widetilde f)\\
   |\gamma|=R}}
 W(\gamma).
 \label{eq:kernel-genealogico-finito-rev2}
\end{equation}
Deux routes ayant le même mot visible restent des termes distincts de la somme
lorsqu’elles appartiennent à des fibres généalogiques différentes.


\begin{proposition}[Composition par coupe enrichie]
\label{prop:composicion-kernel-enriquecido-rev2}
Si tout chemin de longueur \(R+S\) possède une coupure unique après \(R\) transitions, alors
\[
 \mathcal K_{R+S}(\widetilde f,\widetilde i)
 =\sum_{\widetilde x}
 \mathcal K_S(\widetilde f,\widetilde x)
 \mathcal K_R(\widetilde x,\widetilde i).
\]
La formule conserve l’ordre causal des facteurs lorsque \(\mathcal B\) n’est pas
commutative.
\end{proposition}

\begin{proof}
On partitionne l’ensemble des routes complètes selon l’état enrichi
\(\widetilde x\) atteint à la coupure. La factorisation fonctorielle du poids
transforme la somme sur chaque classe en produit des noyaux partiels.
L’unicité de la coupure garantit que chaque route apparaît exactement une fois.

\end{proof}

\subsection{Construction de la phase interne au moyen du caractère réel d'action sur des chemins généalogiques}

Soit \((E,J_E)\) un espace euclidien réel muni d’un endomorphisme
orthogonal \(J_E\) tel que \(J_E^2=-I\). Si
\(\mathcal A_*:\mathsf P_{\rm TPK}\to\mathbb R\) est additive par
concaténation, la phase
\begin{equation}
 \Phi_J(\gamma)
 :=\exp[-J_E\mathcal A_*(\gamma)]
 =\cos\mathcal A_*(\gamma)\,I
  -\sin\mathcal A_*(\gamma)\,J_E
 \label{eq:fase-interna-caminos-rev2}
\end{equation}
est multiplicative. Fixons une section positive non nulle
\(\hbar_{\rm int}\) de la droite d’action \(\mathcal L_S\)
construite précédemment ; on peut utiliser la section
\(\hbar_{\rm ret}\) du profil déclaré.
Pour une action dimensionnelle
\(S_{\rm int}(\gamma)=\hbar_{\rm int}\mathcal A_*(\gamma)\), cette
expression réalise en interne la phase
\(\exp[-J_E\,S_{\rm int}(\gamma)/\hbar_{\rm int}]\). Le symbole scalaire
\(\mathrm i\) apparaît uniquement lors du choix d’un système de coordonnées complexe pour le couple
\((E,J_E)\).

\begin{proposition}[Foncteur de phase]
\label{prop:fase-interna-functor-rev2}
L’application \(\Phi_J\) vérifie
\[
 \Phi_J(\operatorname{id}_x)=I,
 \qquad
 \Phi_J(\gamma_2\circ\gamma_1)
 =\Phi_J(\gamma_2)\Phi_J(\gamma_1).
\]
\end{proposition}

\begin{proof}
L’action du chemin vide est nulle. Pour des routes composables, l’additivité de
\(\mathcal A_*\) et la commutation des multiples d’un même endomorphisme
\(J_E\) permettent de séparer l’exponentielle de la valeur sommée.
\end{proof}

\subsection{Réalisation unitaire sur le registre chronologique
\texorpdfstring{\(27\times27\)}{27 x 27}
}
\label{sec:realizacion-unitaria-registro-27-rev3}
\index{registre chronologique!réalisation unitaire
}
\index{Fourier!blocs du registre chronologique
}
\index{Hamiltoniano!efectivo}

Avant de passer à la limite cylindrique, un noyau fini admet une réalisation
unitaire explicite. Soit

\[
 \mathbb T_{27}^{\,2}:=(\mathbb Z/27\mathbb Z)^2,
 \qquad
 \mathcal D:=\{\mathrm N,\mathrm S,\mathrm E,\mathrm W\},
\]
et considérons l’espace réel
\[
 \mathcal H_{27}^{\mathbb R,\mathrm{reg}}
 :=\ell^2(\mathbb T_{27}^{\,2};\mathbb R)
   \otimes\mathbb R^{\mathcal D}\otimes E.
\]
L'endomorphisme \(J_E\) s'étend comme l'identité sur les facteurs de position et de direction, et la même notation est conservée pour cette extension.

Dans la base ordonnée
\((\mathrm N,\mathrm S,\mathrm E,\mathrm W)\), le mélangeur directionnel est

\begin{equation}
 C_4:=\frac12
 \begin{pmatrix}
  1&1&1&1\\
  1&1&-1&-1\\
  1&-1&1&-1\\
  1&-1&-1&1
 \end{pmatrix},
 \qquad C_4^{\mathsf T}C_4=I_4.
 \label{eq:mezclador-walsh-registro-27-rev3}
\end{equation}
Ici, \(C_4\) désigne exclusivement le mélangeur de Walsh--Hadamard sur la
fibre directionnelle ; lorsqu’il agit sur
\(\mathcal H_{27}^{\mathbb R,\mathrm{reg}}\), on entend
\(C_4\otimes I_E\).

Fixons un système de coordonnées translationnel dans lequel l’action d’une arête élémentaire
\(\gamma_d\) ne dépend que de sa direction :
\[
 s_d:=S_{\rm int}(\gamma_d)\in\mathcal L_S,
 \qquad
 a_d:=\frac{s_d}{\hbar_{\rm int}}\in\mathbb R.
\]
L’opérateur réel diagonal de phase est
\[
 D_J:=\operatorname{diag}
 \left(
  e^{-J_Ea_{\mathrm N}},
  e^{-J_Ea_{\mathrm S}},
  e^{-J_Ea_{\mathrm E}},
  e^{-J_Ea_{\mathrm W}}
 \right).
\]
Le déplacement conditionné transporte chaque composante dans sa direction :
\begin{align*}
 (S\psi)(x,y,\mathrm N)&=\psi(x,y-1,\mathrm N),&
 (S\psi)(x,y,\mathrm S)&=\psi(x,y+1,\mathrm S),\\
 (S\psi)(x,y,\mathrm E)&=\psi(x-1,y,\mathrm E),&
 (S\psi)(x,y,\mathrm W)&=\psi(x+1,y,\mathrm W),
\end{align*}
avec toutes les coordonnées prises modulo \(27\). Le pas élémentaire et sa
composition de douze pas sont
\begin{equation}
 U_{\rm micro}:=S D_J C_4,
 \qquad
 U_{12}^{\mathrm{reg}}:=U_{\rm micro}^{12}.
 \label{eq:paso-unitario-registro-27-rev3}
\end{equation}
L’exposant \(\mathrm{reg}\) conserve le type : cet opérateur agit sur
le tore directionnel \(27\times27\) et ne s’identifie ni aux autres objets
dodécaphasiques de la construction ni à l’endomorphisme temporel d’ordre \(108\).


\begin{proposition}[Unitarité et somme exacte des chemins dans le registre]
\label{prop:unitariedad-suma-registro-27-rev3}
L’opérateur \(U_{\rm micro}\) est orthogonal sur la réalification et unitaire
dans le système de coordonnées complexe déterminé par \(J_E\). En particulier,
\(U_{12}^{\mathrm{reg}}\) est unitaire. Si
\(\Xi:=\mathbb T_{27}^{\,2}\times\mathcal D\), alors, pour tout
\(n\geq0\) et \(\xi_i,\xi_f\in\Xi\),
\[
 \bigl(U_{\rm micro}^{n}\bigr)_{\xi_f,\xi_i}
 =
 \sum_{\substack{\gamma:\xi_i\to\xi_f\\|\gamma|=n}}
 W_{U_{\rm micro}}(\gamma),
\]
où, pour \(\gamma=(\xi_0,\ldots,\xi_n)\),
\[
 W_{U_{\rm micro}}(\gamma)
 :=
 (U_{\rm micro})_{\xi_n,\xi_{n-1}}
 \cdots
 (U_{\rm micro})_{\xi_1,\xi_0}.
\]
La somme parcourt exactement les chemins permis par les blocs non nuls
de l’opérateur.
\end{proposition}

\begin{proof}
Comme \(J_E\) est orthogonal et \(J_E^2=-I\), on a
\(J_E^{\mathsf T}=-J_E\). Par conséquent, chaque bloc
\(e^{-J_Ea_d}\) est une rotation orthogonale et \(D_J\) est orthogonal.
L’opérateur \(S\) permute une base orthonormale et
\eqref{eq:mezclador-walsh-registro-27-rev3} démontre que \(C_4\) est
orthogonal. En conséquence,
\[
 U_{\rm micro}^{\mathsf T}U_{\rm micro}
 =C_4^{\mathsf T}D_J^{\mathsf T}S^{\mathsf T}
   S D_JC_4
 =I.
\]
Les trois facteurs commutent avec l’extension de \(J_E\) ; ainsi,
\(U_{\rm micro}\) est complexe linéaire dans le système de coordonnées
\(J_E\leftrightarrow\mathrm i\), et son orthogonalité réelle équivaut à
l’unitarité. Toute puissance, en particulier
\(U_{12}^{\mathrm{reg}}\), conserve cette propriété.

Pour \(n\geq1\), la multiplication par blocs donne
\[
 \bigl(U_{\rm micro}^{n}\bigr)_{\xi_f,\xi_i}
 =
 \sum_{\xi_1,\ldots,\xi_{n-1}}
 (U_{\rm micro})_{\xi_f,\xi_{n-1}}
 \cdots
 (U_{\rm micro})_{\xi_1,\xi_i}.
\]
Chaque choix d'états intermédiaires détermine un chemin de longueur \(n\), et chaque chemin détermine un unique produit ordonné. Pour \(n=0\), l'unique terme est le chemin vide et son poids est l'identité.
\end{proof}

L’invariance par translations décompose le problème en blocs de quatre
composantes directionnelles, tensorisés avec \(E\). Pour

\[
 k_x=\frac{2\pi a}{27},
 \qquad
 k_y=\frac{2\pi b}{27},
 \qquad 0\leq a,b<27,
\]
nous adoptons la convention de Fourier réalifiée
\[
 \chi_{a,b}(x,y)
 :=\exp[-J_E(k_xx+k_yy)].
\]
La substitution directe dans le décalage conditionné produit
\[
 S(k_x,k_y)
 =\operatorname{diag}
 \bigl(
  e^{J_Ek_y},e^{-J_Ek_y},
  e^{J_Ek_x},e^{-J_Ek_x}
 \bigr),
\]
et, comme \(D_J\) et \(C_4\) ne dépendent pas de la position,
\begin{equation}
 \begin{aligned}
 U_{\rm micro}(k_x,k_y)
   &=S(k_x,k_y)D_JC_4,\\
 U_{12}^{\mathrm{reg}}(k_x,k_y)
   &=U_{\rm micro}(k_x,k_y)^{12}.
 \end{aligned}
 \label{eq:bloques-fourier-registro-27-rev3}
\end{equation}
Par conséquent,
\[
 \operatorname{Spec}U_{12}^{\mathrm{reg}}
 =
 \bigcup_{a,b=0}^{26}
 \operatorname{Spec}
 U_{12}^{\mathrm{reg}}
 \left(\frac{2\pi a}{27},\frac{2\pi b}{27}\right).
\]
La puissance matricielle énumère les chemins et la décomposition de Fourier exprime
les quasi-énergies de la même dynamique finie.

Une durée élémentaire \(\delta t>0\) étant fixée, soit
\(\Delta t_{12}:=12\,\delta t\). Dans chaque bloc unitaire, choisissons une branche
spectrale \(\mathfrak b\) du logarithme. L’opérateur
\(\operatorname{Log}_{\mathfrak b}U_{12}^{\mathrm{reg}}(k_x,k_y)\) est
anti-auto-adjoint et commute avec \(J_E\). On définit
\begin{equation}
 \begin{aligned}
 H_{\rm eff}(k_x,k_y)
 &:=
 \frac{\hbar_{\rm int}}{\Delta t_{12}}\,
 J_E\,
 \operatorname{Log}_{\mathfrak b}
 U_{12}^{\mathrm{reg}}(k_x,k_y),\\
 U_{12}^{\mathrm{reg}}(k_x,k_y)
 &=
 \exp\!\left[
  -\frac{J_E\Delta t_{12}}{\hbar_{\rm int}}\,
  H_{\rm eff}(k_x,k_y)
 \right].
 \end{aligned}
 \label{eq:hamiltoniano-efectivo-registro-27-rev3}
\end{equation}
La première ligne est autoadjointe dans le système de coordonnées complexe. La seconde découle de
\(-J_E^2=I\) et démontre que l’exponentielle retrouve exactement le
propagateur. La branche \(\mathfrak b\) fixe la zone de quasi-énergie.

Cette réalisation construit un propagateur fini, sa somme exacte de chemins et
son générateur effectif. Son identification avec un hamiltonien physique continu
requiert encore de vérifier les hypothèses de représentation, d’action et de limite
énoncées ci-dessous.


\subsection{Limite cylindrique sur l’espace des généalogies}

Les extrémités étant fixées, soient \(\mathcal P_n\) des ensembles finis de chemins à la
résolution \(n\), avec des applications de restriction surjectives. L’espace
\[
 \Omega^{\rm path}_{fi}:=\varprojlim_n\mathcal P_n
\]
est compact et totalement discontinu. Une collection compatible de mesures
de probabilité \(\mu_n\) sur les cylindres détermine une mesure de Borel
\(\mu\) sur \(\Omega^{\rm path}_{fi}\).

Soit \(\Pi_n\) la partition en cylindres de niveau \(n\). Pour
\(C\in\Pi_n\), choisissons un représentant \(\gamma_C\), une action approchée
\(S_n(C)\) et le poids \(\mu_n(C)\). Pour
\(F:\Omega^{\rm path}_{fi}\to\operatorname{End}(E)\), nous définissons

\begin{equation}
 \mathcal I_n(F)
 :=\sum_{C\in\Pi_n}\mu_n(C)F(\gamma_C)
 \exp\!\left[-J_E\frac{S_n(C)}{\hbar_{\rm int}}\right].
 \label{eq:integral-cilindrica-aproximante-rev2}
\end{equation}

\begin{theorem}[Intégrale généalogique de chemins]
\label{thm:integral-genealogica-caminos-rev2}
Supposons que :
\begin{enumerate}
 \item les mesures cylindriques \(\mu_n\) sont projectivement compatibles ;
 \item \(F\) est continue et bornée ;
 \item il existe une action continue \(S:\Omega^{\rm path}_{fi}\to\mathcal L_S\) pour laquelle
 \[
  \max_{C\in\Pi_n}\sup_{\gamma\in C}
  \left\|\frac{S_n(C)-S(\gamma)}{\hbar_{\rm int}}\right\|
  \longrightarrow0;
 \]
 \item le maillage de \(\Pi_n\) tend vers zéro.
\end{enumerate}
Alors \(\mathcal I_n(F)\) converge en norme d’opérateur, la limite est
indépendante des représentants et définit l’intégrale de Bochner
\begin{equation}
 \int_{\Omega^{\rm path}_{fi}}
 F(\gamma)
 \exp\!\left[-J_E\frac{S(\gamma)}{\hbar_{\rm int}}\right]
 \mathrm d\mu(\gamma).
 \label{eq:integral-genealogica-caminos-rev2}
\end{equation}
\end{theorem}

\begin{proof}
La compatibilité des mesures finies construit \(\mu\) sur l'algèbre des cylindres et, par extension, sur la sigma-algèbre de Borel. La fonction
\[
 G(\gamma)=F(\gamma)
 \exp[-J_E\,S(\gamma)/\hbar_{\rm int}]
\]
est continue sur un compact et, par conséquent, uniformément continue et
bornée. L’orthogonalité de \(J_E\) fournit l’estimation

\[
 \|e^{-J_E a}-e^{-J_E b}\|\leq |a-b|.
\]
L’approximation uniforme de l’action et la contraction du maillage font
converger les sommes cylindriques vers l’intégrale de Bochner de \(G\)
\cite{Bochner1955}.
\end{proof}

Le théorème construit une intégrale rigoureuse sur l’espace compact des
généalogies. Sa réalisation comme intégrale oscillante de Feynman pour un
hamiltonien concret s’obtient lorsque les noyaux finis, l’action et la
mesure des cylindres satisfont ces hypothèses de convergence. Le développement
matriciel fini, l’intégrale généalogique et le propagateur physique sont ainsi
reliés par des morphismes explicites, sans être identifiés par abréviation.


\subsection{Valeur exprimée et mémoire de la fibre}

Soit
\(\operatorname{Val}_{\rm path}:\Omega^{\rm path}_{fi}\to\mathcal Y_{\rm val}\subset\mathbb R\)
un lecteur de chemins borélien, où \(\mathcal Y_{\rm val}\) est muni de sa structure
borélienne standard, et soit
\(\nu=(\operatorname{Val}_{\rm path})_*\mu\). Pour toute fonction qui
se factorise par la valeur, \(F=\widehat F\circ\operatorname{Val}_{\rm path}\),
on a
\[
 \int_{\Omega^{\rm path}_{fi}}F\,\mathrm d\mu
 =\int_{\mathcal Y_{\rm val}}\widehat F(x)\,\mathrm d\nu(x).
\]
Comme \(\Omega^{\rm path}_{fi}\) est compact métrisable et \(\mu\) est une
probabilité de Borel, il existe une collection de probabilités conditionnelles
régulières \((\mu_x)_{x\in\mathcal Y_{\rm val}}\), définie pour \(\nu\)-presque tout \(x\). Pour
un intégrande généalogique général, la désintégration de \(\mu\) prend la forme
\begin{equation}
 \int_{\Omega^{\rm path}_{fi}}G(\gamma)\,\mathrm d\mu(\gamma)
 =\int_{\mathcal Y_{\rm val}}\left[
   \int_{\operatorname{Val}_{\rm path}^{-1}(x)}
   G(\gamma)\,\mathrm d\mu_x(\gamma)
  \right]\mathrm d\nu(x).
 \label{eq:desintegracion-genealogica-caminos-rev2}
\end{equation}
L’intégrale intérieure conserve les histoires distinctes qui expriment la même
coordonnée. C’est le lien précis entre la somme sur tous les chemins,
l’évaluation réelle et la thèse selon laquelle le continu est une projection
d’un domaine généalogique plus riche.

\paragraph{Domaine de l’intégrale généalogique de chemins : compacité cylindrique, cohérence projective, convergence des noyaux et désintégration sur les fibres}
Le développement fini, la réalisation unitaire sur
\(\mathbb T_{27}^{\,2}\) et sa réduction de Fourier sont exacts dans leurs domaines
finis. La construction de \(H_{\rm eff}\) est
exacte une fois déclarées la structure complexe, l’action élémentaire
translationnelle, la durée et la branche spectrale ; elle ne sélectionne pas à elle seule un
hamiltonien physique continu. La catégorie enrichie des routes et la phase
interne précèdent l’assemblage cylindrique. Le
\cref{thm:integral-genealogica-caminos-rev2} est exact sous les quatre
prémisses déclarées. La réalisation d’un modèle quantique continu concret
requiert de vérifier ces prémisses pour son action et son hamiltonien.

